\subsection{对应于谢瓦莱系的实形式}

在本节中，我们考虑分裂半单李代数 \((\mathfrak{a},\mathfrak{h})\) （域 \(\mathbf{C}\) 上的；第 VIII 章，§2，第 1 节），其根系为 \(\mathrm{R}(\mathfrak{a},\mathfrak{h}) = \mathrm{R}\) ，以及一个谢瓦莱系 \((X_{\alpha})_{\alpha \in \mathbb{R}}\) （\((\mathfrak{a},\mathfrak{h})\) 的）（第 VIII 章，§2，第 4 节，定义 3）。

回忆（同前），线性映射 \(\theta: a \to a\) ——在 h 上与 \(-Id_{h}\) 相合，且把 \(X_{\alpha}\) 映为 \(X_{-\alpha}\) （对一切 \(\alpha \in R\) ）——是 a 的自同构。此外（同前，命题 7），若 \(\alpha, \beta, \alpha + \beta\) 都是根，则

\[
[ X _ {\alpha}, X _ {\beta} ] = \mathrm{N} _ {\alpha , \beta} X _ {\alpha + \beta}\tag{3}
\]

其中 \(\mathrm{N}_{\alpha ,\beta}\in \mathbf{R}^*\) ，且

\[
\mathrm{N} _ {- \alpha , - \beta} = \mathrm{N} _ {\alpha , \beta}.\tag{4}
\]

用 \(\mathfrak{h}_0\) 表示 \(\mathfrak{h}\) 的由 \(H\in \mathfrak{h}\) （满足 \(\alpha (H)\in \mathbf{R}\) 对一切 \(\alpha \in \mathbb{R}\) ）组成的实向量子空间。则 \(\mathfrak{h}_0\) 是一个 \(\mathbf{R}\) -结构（在复向量空间 \(\mathfrak{h}\) 上），且 \([X_{\alpha},X_{-\alpha}]\in \mathfrak{h}_0\) 对一切 \(\alpha \in \mathbb{R}\) 成立，而 \(\mathfrak{a}\) 的 Killing 型 B 在 \(\mathfrak{h}_0\) 上的限制是分离的正型（第 VIII 章，§2，第 2 节，注 2）。此外，

\[
\mathrm{B} (H, X _ {\alpha}) = 0, \quad \mathrm{B} (X _ {\alpha}, X _ {\beta}) = 0 \text {if} \alpha + \beta \neq 0, \quad \mathrm{B} (X _ {\alpha}, X _ {- \alpha}) <   0\tag{5}
\]

（第 VIII 章，§2，第 2 节，命题 1，及第 4 节，引理 3）。

命题 2. a) 实向量子空间 \(\mathfrak{a}_0 = \mathfrak{h}_0 + \sum_{\alpha \in \mathrm{R}} \mathbf{R}X_\alpha\) （\(\mathfrak{a}\) 的）是 \(\mathfrak{a}\) 的实形式，\(\mathfrak{h}_0\) 是它的嘉当子代数。偶 \((\mathfrak{a}_0, \mathfrak{h}_0)\) 是分裂半单实李代数，\((X_\alpha)\) 是它的谢瓦莱系。

\begin{enumerate}
\def\labelenumi{\alph{enumi})}
\setcounter{enumi}{1}
\tightlist
\item
  设 \(\sigma\) 是 \(\mathfrak{a}\) 相对于 \(\mathfrak{a}_0\) 的共轭。则 \(\sigma \circ \theta = \theta \circ \sigma\) 。\(\sigma \circ \theta\) 的不动点所成的集是一个紧实形式 \(\mathfrak{a}_u\) （\(\mathfrak{a}\) 的），\(i\mathfrak{h}_0\) 是它的嘉当子代数。
\end{enumerate}

部分 a) 由前文立即得出。我们证明 b)。由于 \(\sigma\circ\theta\) 与 \(\theta\circ\sigma\) 是从 a 到 a 的两个在 \(a_{0}\) 上相合的半线性映射，它们相合。而 \(\sigma\circ\theta\) 满足第 1 节的条件 (2)，故它是 a 相对于实形式 \(a_{u}\) 的共轭，该实形式由 \(x\in a\) （满足 \(\sigma\circ\theta(x)=x\) 者）组成（命题 1）。对一切 \(\alpha\in R\) ，令

\[
u _ {\alpha} = X _ {\alpha} + X _ {- \alpha}, v _ {\alpha} = i (X _ {\alpha} - X _ {- \alpha}).\tag{6}
\]

则 R-向量空间 \(a_{u}\) 由 \(ih_{0}\) 、诸 \(u_{\alpha}\) 与诸 \(v_{\alpha}\) 生成。更确切地说，若选取 R 的一个房 C，则

\[
\mathfrak {a} _ {u} = i \mathfrak {h} _ {0} \oplus \bigoplus_ {\alpha \in \mathrm{R} _ {+} (\mathrm{C})} \left(\mathbf {R} u _ {\alpha} + \mathbf {R} v _ {\alpha}\right).\tag{7}
\]

显然 \(i\mathfrak{h}_0\) 是 \(\mathfrak{a}_u\) 的嘉当子代数，余下要证 \(B\) 在 \(\mathfrak{a}_u\) 上的限制是负的。而 \(i\mathfrak{h}_0\) 与诸形如 \(\mathbf{R}u_{\alpha}\oplus \mathbf{R}v_{\alpha}\) 的子空间关于 B 两两正交，此由 (5) 得知；B 在 \(i\mathfrak{h}_0\) 上的限制是负的，且

\[
\mathrm{B} (u _ {\alpha}, u _ {\alpha}) = \mathrm{B} (v _ {\alpha}, v _ {\alpha}) = 2 \mathrm{B} (X _ {\alpha}, X _ {- \alpha}) <   0, \quad \mathrm{B} (u _ {\alpha}, v _ {\alpha}) = 0,\tag{8}
\]

由此即得结论。

注. 采用前面的记号，我们有下列公式：

\[
[ h, u _ {\alpha} ] = - i \alpha (h) v _ {\alpha}, [ h, v _ {\alpha} ] = i \alpha (h) u _ {\alpha}, [ u _ {\alpha}, v _ {\alpha} ] = 2 i H _ {\alpha}, (h \in \mathfrak {h})\tag{9}
\]

\[
[ u _ {\alpha}, u _ {\beta} ] = \mathrm{N} _ {\alpha , \beta} u _ {\alpha + \beta} + \mathrm{N} _ {\alpha , - \beta} u _ {\alpha - \beta}, \quad \alpha \neq \pm \beta ,\tag{10}
\]

\[
[ v _ {\alpha}, v _ {\beta} ] = - \mathrm{N} _ {\alpha , \beta} u _ {\alpha + \beta} + \mathrm{N} _ {\alpha , - \beta} u _ {\alpha - \beta}, \quad \alpha \neq \pm \beta ,\tag{11}
\]

\[
[ u _ {\alpha}, v _ {\beta} ] = \mathrm{N} _ {\alpha , \beta} v _ {\alpha + \beta} - \mathrm{N} _ {\alpha , - \beta} v _ {\alpha - \beta}, \quad \alpha \neq \pm \beta ,\tag{12}
\]

（在后三个公式中，按惯例约定 \(\mathrm{N}_{\gamma, \delta} = 0\) 当 \(\gamma + \delta\) 不是根时）。

注意 \(\sum \mathbf{R}u_{\alpha}\) 是 \(\mathfrak{a}\) 的实子代数，即 \(\mathfrak{a}_0 \cap \mathfrak{a}_u\) 。

设 \(\mathrm{Q}(\mathrm{R})\) 是 R 的根权格（第 VI 章，§1，第 9 节）。回忆：对每个同态 \(\gamma : \mathrm{Q}(\mathrm{R}) \to \mathbf{C}^*\) ，对应着初等自同构 \(f(\gamma)\) （\(\mathfrak{a}\) 的），满足 \(f(\gamma)(h) = h\) 对一切 \(h \in \mathfrak{h}\) ，且 \(f(\gamma)X_{\alpha} = \gamma(\alpha)X_{\alpha}\) （第 VIII 章，§5，第 2 节）。

命题 3. 设 \(\mathfrak{g}\) 是 \(\mathfrak{a}\) 的紧实形式，满足 \(\mathfrak{g} \cap \mathfrak{h} = i\mathfrak{h}_0\) 。则存在同态 \(\gamma: \mathrm{Q}(\mathrm{R}) \to \mathbf{R}_{+}^{*}\) 使得 \(\mathfrak{g} = f(\gamma)(\mathfrak{a}_u)\) 。

设 \(\tau\) 是 \(\mathfrak{a}\) 相对于 \(\mathfrak{g}\) 的共轭。由假设，\(\tau(x) = x\) 对 \(x \in i\mathfrak{h}_0\) 成立，故 \(\tau(x) = -x\) 对 \(x \in \mathfrak{h}_0\) 成立。于是，对一切 \(\alpha \in \mathrm{R}\) 与一切 \(h \in \mathfrak{h}_0\) ，

\[
[ h, \tau (X _ {\alpha}) ] = [ - \tau (h), \tau (X _ {\alpha}) ] = - \tau ([ h, X _ {\alpha} ]) = - \tau (\alpha (h) X _ {\alpha});
\]

由此得出 \([h, \tau(X_{\alpha})] = -\alpha(h)\tau(X_{\alpha})\) 对一切 \(h \in \mathfrak{h}_0\) 成立，从而对一切 \(h \in \mathfrak{h}\) 亦然。故存在 \(c_{\alpha} \in \mathbf{C}^*\) 使得 \(\tau(X_{\alpha}) = c_{\alpha}X_{-\alpha}\) 。由于 \([X_{\alpha}, X_{-\alpha}] \in \mathfrak{h}_0\) ，我们有 \([\tau(X_{\alpha}), \tau(X_{-\alpha})] = -[X_{\alpha}, X_{-\alpha}]\) ，故 \(c_{\alpha}c_{-\alpha} = 1\) ；同样，公式 (3) 与 (4) 给出 \(c_{\alpha + \beta} = c_{\alpha}c_{\beta}\) 当 \(\alpha, \beta, \alpha + \beta\) 都是根时。由第 VI 章，§1，第 6 节，命题 19 的推论 2，存在同态 \(\delta: \mathrm{Q}(\mathrm{R}) \to \mathbf{C}^*\) 使得 \(\delta(\alpha) = c_{\alpha}\) 对一切 \(\alpha \in \mathrm{R}\) 成立。

现在证明每个 \(c_{\alpha}\) 都严格正。事实上，\(c_{\alpha}\mathrm{B}(X_{\alpha},X_{-\alpha})=\mathrm{B}(X_{\alpha},\tau(X_{\alpha}))\) ，由于 \(\mathrm{B}(X_{\alpha},X_{-\alpha})\) 是负的，只需证明对 a 的每个非零元素 z 有 \(\mathrm{B}(z,\tau(z))<0\) ；而 a 的每个元素都可写成 \(x+iy\) ，其中 x 与 y 属于 g，且

\[
\mathrm{B} (x + i y, \tau (x + i y)) = \mathrm{B} (x + i y, x - i y) = \mathrm{B} (x, x) + \mathrm{B} (y, y),
\]

由假设 B 在 g 上的限制是负的且分离的，故所述断言得证。

由此得出同态 \(\delta\) 取值于 \(\mathbf{R}_+^*\) ；于是存在同态 \(\gamma : \mathrm{Q}(\mathrm{R}) \to \mathbf{R}_+^*\) 使得 \(\delta = \gamma^{-2}\) 。此时 \(f(\gamma)^{-1}(\mathfrak{g})\) 是 \(\mathfrak{a}\) 的实形式；对应的共轭为 \(\tau' = f(\gamma)^{-1} \circ \tau \circ f(\gamma)\) 。对一切 \(\alpha \in \mathbb{R}\) ，我们有

\[
\tau^ {\prime} (X _ {\alpha}) = f (\gamma) ^ {- 1} (\tau (c _ {\alpha} ^ {- 1 / 2} X _ {\alpha})) = f (\gamma) ^ {- 1} (c _ {\alpha} ^ {1 / 2} X _ {- \alpha}) = X _ {- \alpha},
\]

且 \(\tau'(h) = \tau(h) = h\) 对 \(h \in i\mathfrak{h}_0\) 成立；由此得出 \(\tau'\) 是相对于 \(\mathfrak{a}_u\) 的共轭，从而 \(f(\gamma)^{-1}(\mathfrak{g}) = \mathfrak{a}_u\) 。

